\section{Conformidade com o enunciado}\label{apx:checklist}

\newcommand{\stpronto}{\textcolor{green!45!black}{\textbf{Pronto}}}
\newcommand{\stparcial}{\textcolor{orange!80!black}{\textbf{Parcial}}}
\newcommand{\stpendente}{\textcolor{red!70!black}{\textbf{Pendente}}}

A Tabela~\ref{tab:checklist} relaciona cada exigência do enunciado~\cite{enunciado} ao local do relatório e ao seu
estado. \stpronto: texto ou dados concluídos nesta versão; \stparcial: método, critérios e formato de relato
fixados, mas a evidência depende dos experimentos; \stpendente: ainda a fazer.

{\small
\begin{longtable}{@{}p{5.4cm}p{3.7cm}p{5.4cm}@{}}
\caption{Exigências do enunciado, local no relatório e estado.}\label{tab:checklist}\\
\toprule
Exigência & Onde & Estado\\
\midrule
\endfirsthead
\toprule
Exigência & Onde & Estado\\
\midrule
\endhead
\midrule
\multicolumn{3}{r}{\emph{continua na próxima página}}\\
\endfoot
\bottomrule
\endlastfoot
Escolha do problema e da base de dados; tipo (classificação) &
Seções~\ref{sec:problema} e~\ref{sec:smd} &
\stpronto: classificação binária por ponto no SMD\\[3pt]
Definições do problema e do conjunto de treinamento &
Seções~\ref{sec:problema} e~\ref{sec:smd} &
\stpronto\\[3pt]
2 a 3 exemplos da base de dados &
Seção~\ref{sec:exemplos} &
\stpronto: 3 observações reais da máquina 1-1\\[3pt]
Valores desejados e atributos &
Seções~\ref{sec:problema} e~\ref{sec:smd} &
\stpronto: rótulo 0/1 e 38 atributos\\[3pt]
Quantidade de amostras &
Seção~\ref{sec:smd}; Apêndice~\ref{apx:maquinas} &
\stpronto\\[3pt]
Preparação e limpeza &
Seção~\ref{sec:prep} &
\stpronto\\[3pt]
Balanceamento de classes &
Seção~\ref{sec:balanceamento} &
\stpronto\\[3pt]
Divisão dos conjuntos &
Seção~\ref{sec:divisao} &
\stpronto\\[3pt]
Métricas de avaliação &
Seção~\ref{sec:metricas} &
\stpronto\\[3pt]
Etapa~I: ao menos 2 modelos de classificação &
Seção~\ref{sec:baselines}; Seção~\ref{sec:resultados} &
\stparcial: protocolo definido, execução pendente; o Hopfield não consta da lista do enunciado (confirmar com o
professor)\\[3pt]
Métricas mínimas: acurácia, precisão, revocação, $F$-scores, AUC-ROC e matriz de confusão &
Seção~\ref{sec:metricas}; Tabelas~\ref{tab:res-s} e~\ref{tab:res-conf} &
\stparcial: formato fixado, valores pendentes\\[3pt]
Etapa~II: classificação por \emph{prompting} (descrição da tarefa, \emph{zero-shot}, \emph{few-shot} com
$\{3,5,10,20\}$) &
Seção~\ref{sec:etapa2}; Apêndice~\ref{apx:prompt} &
\stparcial: \emph{prompt} documentado, execução pendente\\[3pt]
Etapa~II: repetir as métricas; tempo por 100 amostras e custo estimado &
Seção~\ref{sec:etapa2}; Tabelas~\ref{tab:res-llm} e~\ref{tab:registro-llm} &
\stparcial: formato fixado, valores pendentes\\[3pt]
Mesmos dados da Etapa~I; documentar \emph{prompt} e versão do modelo &
Seção~\ref{sec:etapa2}; Apêndice~\ref{apx:prompt} &
\stparcial: \emph{prompt} documentado, versão do modelo pendente\\[3pt]
Comparação crítica (pontos positivos e negativos) &
Seção~\ref{sec:critica} &
\stparcial: critérios e análise \emph{a priori} prontos, evidência empírica pendente\\[3pt]
Conclusões: o que funciona, limites e próximos passos &
Seção~\ref{sec:conclusoes} &
\stparcial: limites e próximos passos prontos; ``o que funciona'' depende dos experimentos\\[3pt]
Equipe de no máximo 3 integrantes &
Folha de rosto; Apêndice~\ref{apx:cronograma} &
\stpendente: nomes e instituição\\[3pt]
Relatório técnico &
Este documento &
\stparcial: estrutura completa, resultados pendentes\\[3pt]
Apresentação remota (02 e 04/12) &
Apêndice~\ref{apx:cronograma} &
\stpendente\\
\end{longtable}
}
